\section{Discussion}
\label{sec:discussion}

\textbf{Why ConvNeXt‑V2 works well on a fine‑grained, non‑linearly‑separable dataset.} Section~\ref{sec:dataset_analysis} showed that frozen ImageNet features are not linearly separable for AvocadoLeaf (near-zero Silhouette score, held-out linear-probe accuracy of only 65.48\%), reflecting substantial visual overlap between disease categories. The architecture comparison in Section~\ref{subsec:arch_compare} is consistent with this picture. Only fine-tuned, non-linear models reach usable accuracy, and among them ConvNeXt‑V2 achieves the best trade-off (92.14\% accuracy, 0.900 macro F1) with the fewest parameters (3.4M) of any competitive architecture, including its predecessor ConvNeXt‑V1 (27.8M). This suggests that, for this task, ConvNeXt‑V2's Global Response Normalisation and design refinements over ConvNeXt‑V1 translate into a more sample- and parameter-efficient decision boundary rather than raw capacity being the deciding factor, since several much larger models (ViT, SwinTransformer) do not clearly outperform it.

\textbf{Few-shot learning narrows but does not close the data-efficiency gap.} The same non-linear-separability finding motivated comparing fine-tuning against metric-based few-shot learning (Section~\ref{subsec:fewshot}). ProtoNet more than doubles fine-tuning's accuracy at every shot level tested (e.g. 57.4\% vs. 28.6\% at 20 shots), confirming that a task-adapted embedding space is far more data-efficient than directly fine-tuning a classifier head on a handful of examples. However, ProtoNet's absolute accuracy remains well below the $>$90\% achieved with the full training set, so few-shot learning should currently be viewed as a rapid-onboarding aid for a low-label adaptation strategy or new orchards rather than a replacement for full-scale data collection. RelationNet's near-chance performance (16--19\%, against a 16.7\% random baseline) is a notable negative result. Unlike ProtoNet's fixed distance metric, RelationNet must learn its comparison function from the same limited episodes, and with only 1,302 training-set images available for episodic training this learnable module likely lacks enough data to converge to a useful metric.

\textbf{The model is reliable at detecting what it doesn't know, but its raw confidence is only moderately calibrated.} The uncertainty analysis (Section~\ref{subsec:uncertainty}) shows two distinct results that carry different practical implications. Out-of-distribution detection is strong for both far OOD (AUC 0.994) and near OOD (AUC 0.985) inputs, meaning the model can reliably flag non-leaf images and, importantly, leaves from other crop species as "not avocado" rather than silently forcing them into one of the six trained classes, a valuable safety property for a farmer-facing tool. The raw ECE of 0.0625 is above the common 0.05 threshold for good calibration, so raw softmax confidence should not be read as a literal probability of correctness out of the box. However, we show in Section~\ref{subsec:uncertainty} that a single fitted temperature parameter ($T=1.274$) brings test-set ECE down to 0.0193 at no cost to accuracy, so this is a substantially mitigated by temperature scaling, limitation, and confidence scores surfaced to end users should simply be the temperature-scaled ones.

\textbf{Interpretability evidence is now quantitative and full-test-set, but still leaves failure modes uncharacterised.} The LIME explanations in Section~\ref{subsec:interpretability} localise to symptomatic regions (scorched lesions for \textit{chaymepla}, powdery patches for \textit{moctrang}) rather than background or leaf-margin artefacts, and the deletion-metric evaluation across all 280 test images (mean confidence drop 0.649, 95\% CI [0.618, 0.678], \ref{app:lime_faithfulness}) confirms this is not an artefact of a favourably chosen example. Hiding the $\sim$20\% of the image LIME attributes each prediction to collapses the model's confidence by nearly two-thirds on average. The drop is nonetheless smaller for misclassified images (0.516) than correct ones (0.660) and is lowest for \textit{repsap} (0.416), whose diffuse, non-contiguous damage pattern is a harder fit for a fixed-size superpixel mask. The confusion-matrix error analysis in Section~\ref{subsec:error_analysis} corroborates this from an independent angle: \textit{repsap} is also the class BoVien most often confuses with \textit{chaymepla} (12.0\% of its test images), so its lower LIME faithfulness and its higher misclassification rate point at the same underlying cause, a visually diffuse symptom pattern, rather than two unrelated weaknesses.

\textbf{A language model is a useful explainer of the CNN's decisions, but not a useful automatic error detector.} Section~\ref{subsec:vlm_explain} showed that a small (2B parameter), free, locally-hosted vision-language model can successfully combine the CNN's prediction with LIME's visual attribution to produce grounded, farmer-readable explanations. This was the original motivation for adding a language model to the system, and it works, as the generated text correctly references both the attributed leaf region and the class-appropriate visual symptom in the examples inspected (Figure~\ref{fig:explain_vlm_showcase}). We separately probed a harder use of the same model that we do \emph{not} report success on (\ref{app:llm} gives the full protocol and results), namely whether the explainer's own \texttt{supports\_cnn\_diagnosis} judgement, without ever overriding the CNN, could serve as a free-standing signal for automatically flagging likely CNN mistakes. It reached a precision of 0.40 but a much lower recall (0.21) than a zero-cost baseline that simply flags the top half of the escalated subset by predictive entropy (precision 0.24, recall 0.79), and every mistake it caught was already caught by that entropy baseline. So, at the model scale available under a no-budget constraint, a language model earns its place in this system specifically as an explanation layer on top of an unchanged CNN + LIME pipeline, not as a substitute for the uncertainty estimation the system already has.

\textbf{Limitations.} Several limitations qualify the conclusions above. First, the dataset, while diverse in imaging conditions, was collected in only three provinces of the Central Highlands, and its class sizes range from 165 to 654 images, so both generalisation to other growing regions and performance on the rarest classes (\textit{repsap}, \textit{gisat}) should be treated cautiously. Second, the expert quality-control review (Section~\ref{sec:dataset_construction}) involved a single expert, so the reported agreement reflects consistency with one reviewer's judgement rather than genuine inter-expert reliability. A second independent expert would be needed to measure that directly. Third, and separately, the annotator pool itself is thin: three annotators means every image is settled by a 2-1 majority vote at best, and the three-way tie-break procedure defined in Section~\ref{sec:dataset_construction} was never exercised in practice, so it remains an untested part of the protocol rather than a validated one. Fourth, the intrinsic difficulty metrics and linear-probe accuracy were computed on a validation set of only 279 images, which inflates the variance of these estimates. Fifth, beyond the GPU inference latency reported in Section~\ref{subsec:arch_compare}, Section~\ref{subsec:realtime_demo} now reports genuine on-device latency on an Android phone (ONNX Runtime, CPU only) and an iPhone (native Core ML), alongside laptop latency on Windows and macOS (both via ONNX Runtime). This substantiates edge-deployment claims for mobile hardware specifically, though a dedicated embedded board such as a Raspberry Pi remains untested, no GPU/NNAPI acceleration was enabled for the Android measurement, and the reported numbers exclude the background-removal preprocessing step (Section~\ref{sec:dataset_construction}), so true end-to-end field latency, especially on Android, would be higher than Table~\ref{tab:realtime_demo} reports. Each platform's numbers also come from a single live capture per class rather than a sustained session, so they should be read as an indicative first measurement rather than a fully characterised latency distribution. Sixth, the real cross-domain evaluation (Section~\ref{subsec:real_cross_domain}) is a genuine same-task, independent-site test, but only at the coarse healthy/diseased granularity, since K-Kotagiri does not provide per-disease labels. It does not establish whether BoVien's specific disease-class predictions (\textit{chaymepla}, \textit{domla}, etc.) generalise across sites, only whether the binary healthy-vs-diseased signal does, and on this coarser task it generalises poorly (60.00\% zero-shot accuracy, down from 98.57\% in-domain). A second, independently collected \emph{6-class} AvocadoLeaf dataset would be needed to test disease-level cross-domain generalisation, and remains future work. Seventh, the annotation guideline assigns a single label per image based on the most severe visible symptom (Section~\ref{sec:dataset_construction}), so leaves with co-occurring symptoms (e.g. mealybug damage alongside an unrelated leaf spot) are reduced to one class. This is a source of label ambiguity that the current single-label taxonomy cannot capture, and it likely contributes to the classes with more diffuse or compound symptom patterns being harder to classify and to explain (Section~\ref{subsec:interpretability}'s per-class LIME faithfulness results). Eighth, the unfrozen-layer depth used during transfer learning is not held perfectly constant across architecture families (Section~\ref{sec:proposed_model}), a source of protocol variation not disentangled by this study. ResNet-18 is the only architecture affected, and it does not rank near the top of Table~\ref{tab:arch_compare}, so this is unlikely to change the selection of ConvNeXt-V2 but has not been verified by an equal-depth re-run. Ninth, the ablation study (Section~\ref{subsec:ablation}) surfaces a discrepancy between the training protocol as described in Section~\ref{sec:proposed_model}, which states that all 15 candidates share the same data-augmentation policy, and the actual checkpoint configuration. The ConvNeXt-V2/BoVien checkpoint reported in Table~\ref{tab:arch_compare} was in fact trained with augmentation disabled, and spot-checking other candidates' saved configurations shows this switch was not applied uniformly across the benchmark. We do not have evidence that this changes which architecture ranks best. The ablation shows augmentation \emph{hurts} BoVien specifically, so if anything this makes the benchmark conservative with respect to BoVien rather than favouring it. We have not re-run the full 15-model benchmark under a corrected, uniform augmentation switch before this submission, since doing so does not change which architecture the paper selects and the compute cost of 15 additional single-seed training runs was better spent, within this study's timeline, on the real cross-domain evaluation of Section~\ref{subsec:real_cross_domain}. A fully protocol-matched re-run remains the direct way to close this gap and is prioritised as future work. Tenth, every architecture in Table~\ref{tab:arch_compare}, including the four ablation configurations, was trained once, with a single random seed. The bootstrap intervals in \ref{app:bootstrap} and \ref{app:bootstrap_ablation} quantify uncertainty from the finite 280-image test set, not from training-run-to-training-run variance. A different seed could shift any single model's accuracy by an amount this study does not measure. This matters most for the top six models in Table~\ref{tab:arch_compare}, which are already statistically tied on test-set noise alone, so seed variance is a second, unmeasured source of uncertainty on top of that tie, not a separate concern that only affects lower-ranked models. Multi-seed re-training of at least the top six candidates is left as future work to quantify this directly.
